\chapter{Le complexe de l'indice de volatilité}\label{ch:s2:the-volatility-index-complex}

Le VIX a terminé l'année 2015 en hausse de 2.4~\% sur son niveau de départ. Une position longue en contrats à terme sur le VIX, tenue à une maturité constante d'environ un mois et reconstituée ici à partir
des règlements-livraisons quotidiens de Cboe, a perdu 37~\% sur la même année. Elle a perdu dans douze des quatorze années civiles de 2013 à septembre 2026, soit 99.99~\% au total. La courbe des contrats à terme
est ascendante la plupart du temps, et une position longue descend le long de celle-ci~: le roll-down rémunère les vendeurs à découvert. Le 5 février 2018, les vendeurs à découvert
découvrent l'autre face de la position. Notre indice à maturité constante a progressé de 95.9~\% en une journée, et un produit inverse sur celui-ci a perdu autant, 95.9~\%. Ce chapitre
construit le complexe sur la trajectoire de variance synthétique du chapitre~1~: un indice de type VIX, ses contrats à terme, des produits longs, inverses et à effet de levier, et les contrats à terme
que ces produits doivent négocier à chaque clôture. La construction est \texttt{firm.vixetp}.

\section{Contrats à terme sur le VIX et leur courbe}

Cboe conçoit le VIX pour mesurer l'anticipation du marché sur la volatilité prospective à 30 jours du S\&P 500. Il le calcule à partir des prix des options SPX et
SPXW, un ensemble de strikes hors de la monnaie pondérés comme dans la réplication d'un swap de variance, sur une échéance proche et une échéance suivante, interpolés à une
maturité constante de 30 jours. Un contrat à terme sur le VIX (livre~1, chapitre~25) se règle à l'échéance sur le cours spécial d'ouverture de l'indice, calculé ce matin-là à partir
des options SPX d'une seule échéance située à 30 jours. Jusqu'à cette date, il se négocie au prix que le marché fixe pour
le niveau de l'indice à l'échéance, et non celui d'aujourd'hui.

\begin{definition}[Pente de la structure par terme du VIX]\label{def:s2:the-volatility-index-complex:slope}
La \emph{pente de la structure par terme du VIX}\index{pente de la structure par terme du VIX} (VIX term-structure slope) est la différence, ou le rapport, entre un contrat à terme sur le VIX plus tardif et un contrat plus proche (souvent le
deuxième et le premier contrat) ou entre le premier contrat et l'indice~; une pente positive (contango) signifie que les positions longues paient un roll-down à mesure que
chaque contrat converge vers l'indice.
\end{definition}

De janvier 2013 à septembre 2026, le deuxième contrat a clôturé au-dessus du premier sur 85.0~\% des jours, de 0.89 points en moyenne, et le premier contrat
se situait 0.63 points au-dessus du VIX. Notre indice long est reconstitué à chaque clôture
entre les deux premiers contrats, avec des pondérations proportionnelles aux jours restants, de sorte que sa maturité reste égale à un cycle de contrat, soit environ un mois.
Une position longue à maturité constante doit sans cesse vendre le contrat le plus proche et acheter le plus éloigné. Lorsque la courbe est
ascendante, chaque contrat à terme glisse vers l'indice à mesure que son échéance approche. Ce glissement est le roll-down de volatilité (livre~5, chapitre~25). Dans notre
reconstitution, l'indice long a perdu à un taux logarithmique de 66.4~\% par an, soit 99.99~\% sur l'ensemble de la période.

\begin{omfigure}
\begin{tikzpicture}
\begin{axis}[width=12cm,height=5.6cm,ybar,bar width=4pt,ylabel={\small \% sur l'année},tick label style={font=\footnotesize},
  xmin=2012.4,xmax=2026.6,ymin=-100,ymax=180,xtick={2013,2015,2017,2019,2021,2023,2025},grid=major,grid style={gray!25},
  table/col sep=comma,/pgf/number format/1000 sep={},legend style={font=\scriptsize,at={(0.5,-0.18)},anchor=north},legend columns=2]
\addplot[fill=omThm!70,draw=omThm] table[x=year,y=long]{figdata/strategies-2/05-the-volatility-index-complex/yearly.csv};
\addplot[fill=omDef!60,draw=omDef] table[x=year,y=vix]{figdata/strategies-2/05-the-volatility-index-complex/yearly.csv};
\legend{position longue à maturité constante en contrats à terme sur le VIX,variation du VIX}
\end{axis}
\end{tikzpicture}

\omcaption{Chaque année civile de 2013 à 2026 (2026 jusqu'au 24 septembre)~: le rendement d'une position longue en contrats à terme sur le VIX, tenue à une maturité constante
d'un cycle de contrats, reconstruit à partir des règlements-livraisons quotidiens de Cboe, et la variation du VIX. La position longue a perdu chaque année sauf en 2018 et 2020. Établi à partir des historiques
de contrats à terme et d'indices de Cboe~; les données brutes ne sont pas redistribuées. Données~: \texttt{s2\_fetch\_vx}.}\label{fig:s2:the-volatility-index-complex:yearly}
\end{omfigure}

L'indice est une illustration année par année de ce que fait le roll-down (\cref{fig:s2:the-volatility-index-complex:yearly}). En 2014, le VIX a progressé de 35~\% et la
position longue a perdu 25~\%. En 2025, le VIX a reculé de 17~\% et la position longue a perdu 45~\%. Seules en 2018 et 2020, années où un pic suffisamment fort dépassait le
roll, la position longue a dégagé un gain.

\subsection*{Le complexe synthétique}

\texttt{firm.vixetp} construit les mêmes objets sur la trajectoire de variance du chapitre~1, avec une volatilité de la variance plus calme (0.25 au lieu de 0.5). L'indice est la
volatilité implicite à 21 jours, soit la variance attendue multipliée par la prime, en points. Les contrats à terme expirent tous les 21 jours de négociation. Chacun est évalué au niveau de l'indice
attendu à l'échéance, revenant vers sa moyenne à la vitesse de la variance, multiplié par $1 + 0.5h$ pour $h$ années jusqu'à l'échéance~: une prime de terme que paient les positions longues
(\cref{lst:s2:the-volatility-index-complex:curve}). L'indice affichait une moyenne de 19.6, avec une médiane de 17.7. Le deuxième contrat se situait au-dessus du premier sur
81~\% des jours, et le premier au-dessus de l'indice de 0.44 points. L'indice long à maturité constante a perdu à un taux logarithmique de 74~\% par an, avec une volatilité de
88~\%. Sans la prime de terme, la courbe reste ascendante sur 65~\% des jours, car l'indice passe l'essentiel de son temps sous sa moyenne. L'indice long
perd ensuite à un taux logarithmique de 26~\% par an, alors que son rendement quotidien moyen est proche de zéro ($-3.7~\%$ par
an)~: la perte vient de la capitalisation d'une série aussi volatile.

\section{Produits cotés de volatilité et leur rebalancement}

\begin{definition}[Produit coté de volatilité]\label{def:s2:the-volatility-index-complex:etp}
Un \emph{produit coté de volatilité}\index{produit coté de volatilité} (exchange-traded volatility product, ETP) est un fonds indiciel coté (ETF) ou un titre de créance coté (livre~3, chapitre~19) qui
délivre un multiple ($1$, $-1$, $2$, $-0.5$) du rendement quotidien d'un indice de contrats à terme sur le VIX, généralement à maturité constante~; il détient ou couvre les
contrats à terme et les rebalance à chaque clôture.
\end{definition}

Un produit qui promet $L$ fois le rendement quotidien, de valeur $A$, détient $LA$ en contrats à terme. Après un rendement quotidien $r$, la valeur est $A(1 + Lr)$ et les
contrats détenus valent $LA(1 + r)$. Pour détenir $L$ fois la nouvelle valeur, il doit acheter
\[
LA(1 + Lr) - LA(1 + r) = A\,L(L - 1)\,r.
\]

\begin{definition}[Flux de rebalancement des ETP]\label{def:s2:the-volatility-index-complex:flow}
Le \emph{flux de rebalancement des ETP}\index{flux de rebalancement des ETP} (ETP rebalancing flow) est la quantité de contrats à terme que les produits de volatilité à effet de levier et inverses doivent négocier à
la clôture pour rétablir leurs multiples, $A\,L(L - 1)\,r$ pour chaque produit~; il est nul pour $L = 1$ et a le signe du mouvement de la journée pour $L = -1$ et
$L = 2$, de sorte que ces produits achètent après les hausses et vendent après les baisses.
\end{definition}

Pour un produit inverse, $L(L - 1) = 2$. Un jour où l'indice de contrats à terme progresse de 50~\% le force à acheter des contrats valant l'ensemble de sa valeur à la clôture, et
les positions longues à double effet de levier doivent faire de même. Le flux est connu à l'avance, sous la forme d'une formule du mouvement de la journée, et il arrive à la clôture lorsque la liquidité est
la plus faible. L'avertissement de Whaley est venu tôt~: les produits sur le VIX les plus populaires ne conviennent pas à des placements de type «~acheter et conserver~» et sont pratiquement assurés de perdre
de l'argent avec le temps. Les produits fondés sur les indices de contrats à terme à court terme avaient perdu près de 4~milliards de dollars depuis leur lancement en 2009.

\section{Contrats à terme contre options}

Un contrat à terme sur le VIX et les options SPX évaluent des grandeurs liées. Les options évaluent la variance du S\&P 500 sur une fenêtre, et le contrat à terme évalue la racine carrée de
une variance à 30 jours qui débute à son échéance. La racine carrée est concave, de sorte qu'un contrat à terme sur la volatilité vaut moins que la racine carrée de la
variance forward~: l'écart est un ajustement de convexité, qui croît avec la volatilité de la volatilité. Un trader qui juge que l'ajustement dans les contrats à terme
est erroné peut arbitrer le contrat à terme contre une variance forward construite à partir d'options SPX à deux échéances. Ce qu'il détient, c'est la volatilité de
la volatilité, et les options sur le VIX (livre~5, chapitre~14) la valorisent directement. Lorsque l'indice bondit, les trois bougent ensemble. Le risque du book de relative value est
qu'ils évoluent alors de façon différente~: le contrat à terme se règle sur le cours spécial d'ouverture, les options sur leur propre cours.

\section{Février 2018}

Le complexe a été testé le 5 février 2018. Le VIX a clôturé à 17.31 le vendredi 2 février et à 37.32 le lundi 5 février. Le premier contrat à terme a clôturé à
15.625, puis à 33.225. Notre indice à maturité constante a progressé de 95.9~\% le lundi et a reculé de 25.8~\% le mardi, ses meilleures et pires journées de 2013 à 2026. Un
produit inverse sur celui-ci avait perdu 96.7~\% depuis le début de l'année.

Credit Suisse a annoncé le 6 février que son titre de créance inverse sur le VIX, XIV, avait subi un événement d'accélération~: sa valeur indicative intrajournalière du 5 février
avait été inférieure ou égale à 20~\% de la valeur indicative de clôture de la veille, 108.3681~\$. Le titre devait être remboursé à sa valeur à la date d'évaluation accélérée
Augustin, Cheng et Van den Bergen ont décrit l'épisode comme une leçon sur le rebalancement de la couverture et de l'effet de levier lorsque les marchés sont concentrés
et volatils, et l'ont comparé à l'assurance de portefeuille. Sur le mouvement de 95.9~\% de notre reconstitution, le flux de rebalancement d'un produit inverse vaut 1.92 fois sa valeur, le tout
acheté à la clôture.

Le complexe synthétique a son propre épisode de février. Après six années de calme, le premier krach simulé fait passer l'indice de 17.3 à 63.6 en une journée, et la position longue
à maturité constante progresse de 189~\%. Un produit inverse assorti de la clause d'accélération avait atteint 16.5 fois sa valeur de départ~; il finit cette journée à
zéro. À cette clôture, il doit acheter des contrats valant 3.77 fois sa valeur précédente, soit deux fois le mouvement de la journée, et une position longue à double effet de levier doit en faire autant
(\cref{fig:s2:the-volatility-index-complex:products}). Un book vendeur de contrats à terme représente la même opération, à une taille choisie~:

\begin{center}
\small
\begin{tabular}{@{}p{5.2cm}>{\centering\arraybackslash}p{2.3cm}>{\centering\arraybackslash}p{2.3cm}>{\centering\arraybackslash}p{2.3cm}>{\centering\arraybackslash}p{2.3cm}@{}}
\toprule
Vente de contrats à terme, avec rebalancement quotidien & rendement annuel & drawdown maximal & valeur finale (départ 1) & survie\\
\midrule
0.10 du capital & 4.8~\% & $-27~\%$ & 2.57 & oui\\
0.25 du capital & 9.6~\% & $-67~\%$ & 6.26 & oui\\
0.50 du capital & --- & $-100~\%$ & 0 & non\\
1.00 du capital & --- & $-100~\%$ & 0 & non\\
\bottomrule
\end{tabular}
\end{center}

\begin{omfigure}
\begin{tikzpicture}
\begin{axis}[width=11cm,height=5.6cm,ymode=log,xlabel={\small année},ylabel={\small valeur (départ 1)},tick label style={font=\footnotesize},
  xmin=0,xmax=20,ymin=0.5,ymax=30,grid=major,grid style={gray!25},table/col sep=comma,unbounded coords=jump,
  legend style={font=\scriptsize,at={(0.5,-0.3)},anchor=north},legend columns=2]
\addplot[omThm,thick] table[x=year,y=inverse]{figdata/strategies-2/05-the-volatility-index-complex/products.csv};
\addplot[omDef,thick] table[x=year,y=quarter]{figdata/strategies-2/05-the-volatility-index-complex/products.csv};
\legend{produit inverse (s'annule au premier krach),vente de contrats à terme à 0.25 du capital}
\end{axis}
\end{tikzpicture}

\omcaption{Deux positions courtes sur l'indice synthétique de contrats à terme sur le VIX à maturité constante, en échelle logarithmique~: un produit inverse assorti d'une clause d'accélération à
20~\% de la clôture précédente, qui atteint 16.5 fois sa valeur et s'annule au premier krach simulé (année 6.0), et un book vendeur à un quart de son
capital, qui survit à trois krachs. Données~: \texttt{s2\_vixetp.paths}.}\label{fig:s2:the-volatility-index-complex:products}
\end{omfigure}

La taille est toute la stratégie. À un quart de son capital, le book gagne 9.6~\% par an et survit à un drawdown de 67~\%. À moitié, il est anéanti. Le
roll-down est réel, mais son ampleur ne dit rien de la taille qui survit le jour où l'indice triple.

\section{Fiches de stratégie}

\begin{strategyfile}{Roll-down de la vente de contrats à terme sur le VIX}\label{strat:s2:the-volatility-index-complex:rolldown}
\sfield{Qui vous paie, et pourquoi} Les acheteurs de volatilité, directement ou par des produits longs, qui paient une prime de terme pour une protection qui est
disponible lorsqu'ils en ont besoin.
\sfield{Instruments et plateformes} Contrats à terme sur le VIX, à la Cboe Futures Exchange~; produits inverses.
\sfield{Signal} La pente de la courbe~: position courte lorsqu'elle est en contango, position plus petite ou nulle lorsqu'elle est inversée.
\sfield{Dimensionnement et exécution} Une petite fraction du capital, fixée par un pic qui triple l'indice~; rebalancement quotidien de la maturité.
\sfield{Coûts} Spreads sur contrats à terme~; coûts de roll.
\sfield{Comment elle meurt} Un pic~: le 5 février 2018 a mis fin à un titre inverse dont la valeur a chuté de plus de 80~\% en une journée.
\sfield{Horizon, capacité, infrastructure} Des mois à des années~; accès aux contrats à terme et roll quotidien.
\sfield{Backtest honnête} Inclure 2018 et 2020~; utiliser les règlements-livraisons et non les niveaux d'indice~; dimensionner sur le pire jour et non sur le ratio de Sharpe.
\sfield{Sources} Whaley (2013)~; Credit Suisse (2018)~; Augustin, Cheng et Van den Bergen (2021)~; ce chapitre~: une position longue à maturité constante
a perdu 99.99~\% de 2013 à 2026, et l'inverse a perdu 96.7~\% du 2 janvier au 5 février 2018.
\end{strategyfile}

\begin{strategyfile}{Anticipation du rebalancement des ETP}\label{strat:s2:the-volatility-index-complex:flow}
\sfield{Qui vous paie, et pourquoi} Les produits à effet de levier et inverses dont le rebalancement, $A\,L(L - 1)\,r$, est une formule du mouvement du jour et se négocie à la
clôture.
\sfield{Instruments et plateformes} Contrats à terme sur le VIX avant et à la clôture.
\sfield{Signal} Les valeurs et multiples publiés des produits, et le mouvement du jour jusqu'ici.
\sfield{Dimensionnement et exécution} Petit face au flux~; positions prises avant la clôture et dénouées après.
\sfield{Coûts} Spreads à la clôture, où la liquidité est faible.
\sfield{Comment elle meurt} Les produits rétrécissent ou ferment~; d'autres anticipent le même flux. Négocier avant un flux connu est légal lorsqu'il repose sur une
information publique~; négocier sur l'ordre d'un client ne l'est pas (chapitre~29).
\sfield{Horizon, capacité, infrastructure} Des heures~; une estimation en direct des flux des produits.
\sfield{Backtest honnête} Les exécutions aux cours de clôture sont optimistes les jours qui comptent.
\sfield{Sources} Augustin, Cheng et Van den Bergen (2021) sur le risque de rebalancement~; aucune performance vérifiée.
\end{strategyfile}

\begin{strategyfile}{Contrats à terme sur le VIX contre options sur le SPX}\label{strat:s2:the-volatility-index-complex:options}
\sfield{Qui vous paie, et pourquoi} Mauvaise tarification de l'ajustement de convexité entre le contrat à terme et la variance forward issue des options sur le SPX.
\sfield{Instruments et plateformes} Contrats à terme sur le VIX~; options sur le SPX à deux échéances~; options sur le VIX pour la volatilité de la volatilité.
\sfield{Signal} Le contrat à terme comparé à la racine carrée de la variance forward, ajustée de la convexité.
\sfield{Dimensionnement et exécution} Jambes calibrées en vega~; testées sous contrainte lors d'un pic.
\sfield{Coûts} Spreads des deux marchés~; écarts de règlement-livraison.
\sfield{Comment elle meurt} Un pic où les jambes bougent de montants différents~; règlement-livraison à des prix différents.
\sfield{Horizon, capacité, infrastructure} Des semaines~; surfaces SPX et contrats à terme sur le VIX.
\sfield{Backtest honnête} Modéliser le cours spécial d'ouverture (SOQ) à l'échéance.
\sfield{Sources} Aucune performance vérifiée.
\end{strategyfile}

\begin{strategyfile}{Volatilité longue sur une courbe inversée}\label{strat:s2:the-volatility-index-complex:inverted}
\sfield{Qui vous paie, et pourquoi} Personne ne paie, en moyenne~: c'est une assurance qui coûte peu lorsque la courbe est inversée, car le roll joue alors en
faveur des positions longues.
\sfield{Instruments et plateformes} Contrats à terme sur le VIX ou options d'achat.
\sfield{Signal} Une courbe inversée (premier contrat au-dessus du second).
\sfield{Dimensionnement et exécution} Petit~; maintenue tant que la courbe est inversée.
\sfield{Coûts} Spreads~; plus larges lors des pics.
\sfield{Comment elle meurt} La courbe revient rapidement au contango et le roll se retourne contre la position longue.
\sfield{Horizon, capacité, infrastructure} De quelques jours à quelques semaines.
\sfield{Backtest honnête} Les inversions sont rares (15~\% des jours de 2013 à 2026)~; les résultats reposent sur quelques épisodes.
\sfield{Sources} Statistiques de courbe de ce chapitre~; aucune performance vérifiée.
\end{strategyfile}

\section{Tutoriel~: roll-down}\label{tut:s2:the-volatility-index-complex:tut}

\begin{tutorial}
\textbf{Objectif.} Construire l'indice synthétique, ses contrats à terme et ses produits~; mesurer le roll-down, le flux de rebalancement et le premier krach~; reconstruire l'indice réel
à maturité constante à partir des règlements-livraisons de Cboe. \textbf{État final~:} le tableau et les deux figures.
\begin{tutsteps}
\item \textbf{Contrats à terme et indice à maturité constante}.
\omcode{code/firm/vixetp/firm_vixetp.py}{36}{62}{Contrats à terme avec prime de terme et indice long à maturité constante.}\label{lst:s2:the-volatility-index-complex:curve}
\item \textbf{Produits et leurs flux}.
\omcode{code/firm/vixetp/firm_vixetp.py}{65}{82}{Produit à multiple quotidien doté d'une clause d'accélération, et son flux de rebalancement.}\label{lst:s2:the-volatility-index-complex:etp}
\item \textbf{Exécuter} \texttt{s2\_fetch\_vx.py} une fois (environ 170 fichiers~; définir \texttt{OQB\_CACHE}), puis \texttt{curve\_stats()}, \texttt{short\_book()},
\texttt{spike()} et \texttt{fig\_vixetp.py}.
\end{tutsteps}
\textbf{Pour aller plus loin.} Dimensionner le book vendeur selon la pente de la courbe~; ajouter une couverture par un call sur le VIX et la chiffrer~; donner des actifs aux produits et
des positions ouvertes aux contrats à terme, puis mesurer le flux en proportion de celles-ci.
\end{tutorial}

\section{Construction~: indice de volatilité et produits}\label{bld:s2:the-volatility-index-complex:vixetp}

\begin{build}
\bfield{Rôle} Un indice de type VIX, ses contrats à terme, un indice à maturité constante, des produits à multiple quotidien et leurs flux de rebalancement.
\bfield{Interface} \texttt{vix\_index(v, cfg)}, \texttt{future(x, mean, cfg, h, premium)}, \texttt{curve\_path(v, cfg, premium, every)},
\texttt{etp(r, leverage, accel)}, \texttt{flow(product, r, leverage)}.
\bfield{Règles} Le contrat arrivant à échéance se règle sur l'indice~; les pondérations sont fixées à la clôture précédente~; un produit s'annule lorsqu'une journée le laisse à un niveau égal ou inférieur à la
fraction d'accélération de la clôture précédente.
\bfield{Tests d'acceptation} \texttt{code/firm/vixetp/tests/}~: la formule du flux à la main pour $L = -1, 1, 2$~; un produit éteint, figé sans
flux ultérieur~; un indice plat sans prime donne une courbe plate et aucun roll, et avec une prime un contango et un roll-down régulier.
\bfield{Pour aller plus loin} Un second facteur de volatilité, plus rapide~; options sur le VIX portant sur l'indice~; actifs des produits et positions ouvertes.
\end{build}

\begin{omsources}
\item Cboe, Volatility Index Methodology~: Cboe Volatility Index.
\item R.~E.~Whaley, «~Trading volatility~: at what cost~?~», \emph{Journal of Portfolio Management} 40(1), 2013.
\item P.~Augustin, I.-H.~Cheng et L.~Van den Bergen, «~Volmageddon and the failure of short volatility products~», \emph{Financial Analysts Journal}
77(3), 2021.
\item Credit Suisse AG, «~Credit Suisse AG announces the event acceleration of its XIV ETNs~», communiqué de presse, 6 février 2018 (SEC exhibit 99.1).
\item Cboe Futures Exchange, VIX futures daily settlements~; Cboe, VIX history.
\end{omsources}

\section{Exercices}

\begin{exercise}[$\star$]\label{exo:s2:the-volatility-index-complex:1}
Un produit rapporte $-1$ fois l'indice chaque jour et vaut 100. L'indice monte de 30~\%. Que doit-il acheter à la clôture~?
\end{exercise}

\begin{exercise}[$\star$]\label{exo:s2:the-volatility-index-complex:2}
Même question pour un produit dont le rendement est égal à deux fois l'indice, et pour un produit dont le rendement est égal à l'indice.
\end{exercise}

\begin{exercise}[$\star$]\label{exo:s2:the-volatility-index-complex:3}
Le premier contrat a 10 jours avant échéance, le second 38, et les contrats sont espacés de 28 jours. Quel poids une position à maturité constante de 30 jours accorde-t-elle au premier~?
\end{exercise}

\begin{exercise}[$\star\star$]\label{exo:s2:the-volatility-index-complex:4}
En 2014, le VIX a progressé de 35~\% et la position longue à maturité constante a perdu 25~\%. Expliquer.
\end{exercise}

\begin{exercise}[$\star\star$]\label{exo:s2:the-volatility-index-complex:5}
Pourquoi un contrat à terme sur la volatilité vaut-il moins que la racine carrée de la variance forward~?
\end{exercise}

\begin{exercise}[$\star\star$]\label{exo:s2:the-volatility-index-complex:6}
Pourquoi les \omterm{def:s2:the-volatility-index-complex:etp}{produits cotés de volatilité} inverses et à effet de levier aggravent-ils un pic à la clôture~?
\end{exercise}

\begin{exercise}[$\star\star\star$]\label{exo:s2:the-volatility-index-complex:7}
\emph{Programmation.} Exécuter \texttt{curve\_stats(0.0)}. Sans prime de terme, pourquoi la courbe reste-t-elle orientée à la hausse la plupart des jours, et pourquoi l'indice long
perd encore~?
\end{exercise}

\begin{exercise}[$\star\star\star$]\label{exo:s2:the-volatility-index-complex:8}
\emph{Trouver la faille.} «~Les positions vendeuses sur contrats à terme sur le VIX ont rapporté 9.6~\% par an avec un quart du capital, donc avec la totalité du capital elles rapporteraient près de 40~\%.~»
\end{exercise}

\section{Problème~: roll-down de volatilité}

\begin{problem}[{Problème du week-end --- volatilité courte via des contrats à terme}]\label{pb:s2:the-volatility-index-complex:1}
Le complexe synthétique, les contrats à terme de Cboe et février 2018.

\medskip\noindent\textbf{Partie I --- La courbe.}
\begin{enumerate}
\item Que mesure le VIX, et comment est-il calculé~?
\item Définir la \omterm{def:s2:the-volatility-index-complex:slope}{pente de la structure par terme du VIX} et donner ses statistiques réelles.
\item Qu'est-ce que le roll-down, et combien a-t-il coûté à une position longue à maturité constante~?
\item Combien la position longue a-t-elle perdu en 2014, 2015 et 2025, et qu'a fait le VIX~?
\end{enumerate}

\medskip\noindent\textbf{Partie II --- Produits.}
\begin{enumerate}[resume]
\item Définir un \omterm{def:s2:the-volatility-index-complex:etp}{produit coté de volatilité}.
\item Dériver le flux de rebalancement $A\,L(L - 1)\,r$.
\item Quels produits achètent après des hausses~?
\item Qu'a trouvé Whaley~?
\end{enumerate}

\medskip\noindent\textbf{Partie III --- Février 2018.}
\begin{enumerate}[resume]
\item Que s'est-il passé pour le VIX et le premier contrat à terme le 5 février~?
\item Qu'ont fait l'indice à maturité constante et son inverse~?
\item Qu'a annoncé Credit Suisse, et pourquoi~?
\item Comment Augustin, Cheng et Van den Bergen ont-ils interprété l'épisode~?
\end{enumerate}

\medskip\noindent\textbf{Partie IV --- Le verdict.}
\begin{enumerate}[resume]
\item Énoncer le \emph{résultat nommé}~: le rendement du roll-down du portefeuille vendeur de contrats à terme et la perte du produit inverse lors du pic artificiel.
\item Combien les portefeuilles vendeurs ont-ils gagné pour chaque taille~?
\item Quel flux le produit inverse a-t-il dû affronter à la clôture du krach~?
\item Pourquoi la prime de terme n'est-elle pas nécessaire pour que la courbe soit orientée à la hausse~?
\item Comment on dimensionnerait un portefeuille vendeur de contrats à terme sur le VIX~?
\item Comment on backtesterait cela honnêtement~?
\item Quelle fiche de stratégie tire profit du rebalancement des autres~?
\item En une phrase~: qui paie le roll-down, et pourquoi~?
\end{enumerate}
\end{problem}

\section{Questions d'entretien}

\begin{interviewq}[$\star$ \iqroles{trader}]\label{iq:s2:the-volatility-index-complex:1}
Pourquoi les produits longs sur le VIX perdent-ils de l'argent avec le temps~?
\end{interviewq}

\begin{interviewq}[$\star\star$ \iqroles{trader}]\label{iq:s2:the-volatility-index-complex:2}
La courbe du VIX est inversée. Que vous indique-t-elle, et quelles transactions suggère-t-elle~?
\end{interviewq}

\begin{interviewq}[$\star\star$ \iqroles{chercheur}]\label{iq:s2:the-volatility-index-complex:3}
Comment construiriez-vous un indice à maturité constante à partir des cours de règlement-livraison des contrats à terme, et qu'est-ce qui peut mal tourner au roll~?
\end{interviewq}

\begin{interviewq}[$\star\star$ \iqroles{risque}]\label{iq:s2:the-volatility-index-complex:4}
Vous gérez un portefeuille vendeur de contrats à terme sur le VIX. Comment le dimensionnez-vous et le soumettez-vous à des tests de résistance~?
\end{interviewq}

\begin{interviewq}[$\star\star$ \iqroles{développeur}]\label{iq:s2:the-volatility-index-complex:5}
Concevez le processus quotidien qui calcule le rebalancement d'un produit à effet de levier à la clôture.
\end{interviewq}

\begin{interviewq}[$\star\star\star$ \iqroles{chercheur}]\label{iq:s2:the-volatility-index-complex:6}
Montrez qu'un produit qui rapporte $L$ fois le rendement quotidien doit négocier $A\,L(L - 1)\,r$ à chaque clôture, et montrez que sa valeur sur deux jours n'est pas $L$
fois le rendement de l'indice sur deux jours.
\end{interviewq}
